%!TEX root = ../../main.tex
%% ===============================================================
%% 10.2 설계 결정에 대한 논의
%% 파일: chapters/10_discussion/02_design_decisions.tex
%% 상위: chapters/10_discussion/chapter.tex
%% 정의 라벨: sec:disc-design
%% 참조 라벨: -
%% 포함 객체: -
%% 인용: -
%% 상태: 초안 (docs/OUTLINE.md 참고)
%% ===============================================================

\section{설계 결정에 대한 논의}
\label{sec:disc-design}

\paragraph{왜 킬만으로 교전을 정의하는가.} 근접성만으로 교전을 정의하면 서로 가까이 있으나 싸우지 않는 상황(파밍, 대치)이 대량으로 섞인다. 킬은 공개 텔레메트리에서 전투가 실제로 일어났음을 보장하는 유일한 사건이다. 대가는 킬 없는 교전의 누락이며, 이는 정의의 범위 제한으로 명시하고 검증 프로토콜에서 별도로 계량한다.

\paragraph{왜 위치를 선택적으로 0으로 두는가.} 지도 위치는 그래프 모델에 필수적인 구조 정보이지만, 순차 모델에서는 ``바론 둥지 근처의 교전은 특정 팀이 이긴다''와 같은 위치 편향을 암기하게 만든다. 따라서 $X^{\mathrm{node}}$에는 보존하고 거시 시퀀스에서는 0으로 둔다.

\paragraph{왜 스칼라 상태를 보간하지 않는가.} 아이템 구매, 버프 만료, 레벨 상승은 이산 전이이다. 프레임 사이를 선형 보간하면 실제로 존재하지 않는 매끄러운 중간 상태가 생기고, 특히 온셋 이후 프레임을 참조하는 보간은 교전 결과 자체를 입력에 새게 한다. 계단 유지는 이 두 문제를 동시에 피한다.

\paragraph{선택 편향의 회피.} 교전의 존재는 킬로 조건화하지만, 라벨 창의 \emph{결과}에는 조건화하지 않는다. 예측 지평에 특정 신호가 있는 창만 남기면 $P(Y \mid X, \text{signal}) \ne P(Y \mid X)$가 되어 Berkson 형 편향이 생기므로, 조용한 라벨 창도 그대로 둔다.
